
The experimental design for validating the four sub-hypotheses was fully specified, though execution was blocked at h-e1.

\subsubsection{Research Questions}

\textbf{RQ1 (Existence)}: Can static analyzers extract SMT constraints from LLM-generated typed code at $\geq 90$\% success rate? (Untested due to infrastructure failure)

\textbf{RQ2 (Mechanism)}: Does dependency analysis accurately identify invalidation cones while maintaining soundness? (Untested, prerequisite RQ1 failed)

\textbf{RQ3 (Performance)}: Does incremental SMT achieve speedup vs batch re-verification on iterative LLM code repair? (Untested, prerequisite RQ1, RQ2 failed)

\subsubsection{Dataset}

\textbf{Source}: HumanEval (OpenAI), problems 0--99 (100 programs).

\textbf{Extension}: Mechanically generated Pydantic type annotations via template-based transformation. Each extended prompt includes function signature with type hints, Pydantic BaseModel schemas for inputs/outputs, @validator decorators encoding preconditions, and original docstring and test suite.

\textbf{Rationale}: HumanEval provides standard baseline for code generation research. Type extension enables static analysis without manual annotation bias. Programs span 10--50 LOC (suitable for proof-of-concept validation).

\textbf{Validation}: Dataset extension pipeline executed successfully, generating 100 JSON files in \texttt{src/data/humaneval\textbackslash \_pydantic/}. Each file contains valid Pydantic template matching the original HumanEval problem specification.

\subsubsection{LLM Code Generation}

\textbf{Model}: Claude Sonnet 3.5 (claude-sonnet-3-5-20240620)

\textbf{Temperature}: 0.2

\textbf{Max Tokens}: 512

\textbf{System Prompt}: ''Generate typed Python code using Pydantic BaseModel. Include @validator decorators for preconditions. Do not use eval, exec, or metaprogramming.''

\textbf{Expected Output}: 100 typed Python programs with Pydantic annotations, passing original HumanEval test suites.

\textbf{Actual Result}: 0 programs generated. All 48 generation attempts failed with HTTP 401 ''API key is invalid'' errors.

\subsubsection{Baseline Methods}

\textbf{Baseline 1: Batch SMT Verification}
\begin{itemize}
\item \textbf{Method}: Re-verify entire program on each repair iteration using Z3.
\item \textbf{Correctness}: Sound and complete.
\item \textbf{Use}: Primary performance baseline for h-m3 speedup measurement (not tested).
\end{itemize}

\textbf{Baseline 2: Test-Suite-Only Validation}
\begin{itemize}
\item \textbf{Method}: LLM generates code, run HumanEval test suite, accept if all tests pass.
\item \textbf{Correctness}: Incomplete (misses edge cases not covered by tests).
\item \textbf{Use}: Comparison point for verification overhead vs no verification (comparison not performed).
\end{itemize}

No baseline comparisons were conducted due to infrastructure failure.

\subsubsection{Evaluation Metrics}

\textbf{h-e1: Extraction Success Rate}
$$\text{extraction\_rate} = \frac{\text{programs with constraints}}{\text{total programs}} \times 100\%$$

\textbf{h-m3: Speedup Factor (Not Measured)}
$$\text{speedup} = \frac{\text{batch verification time}}{\text{incremental verification time}}$$

\textbf{h-m2: Invalidation Cone Size (Not Measured)}
$$\text{cone\_size} = \frac{\text{constraints in invalidation cone}}{\text{total constraints}} \times 100\%$$

\textbf{h-m2: Soundness (Error Detection Rate) (Not Measured)}
$$\text{detection\_rate} = \frac{\text{seeded errors caught}}{\text{total seeded errors}} \times 100\%$$

\textbf{h-m1: Extraction Overhead (Not Measured)}
$$\text{overhead\_ratio} = \frac{\text{extraction time}}{\text{total verification time}} \times 100\%$$

\subsubsection{Threats to Validity}

\textbf{Internal Validity}: Infrastructure failure (invalid API key) blocked testing of core hypothesis. Mitigated by isolating failure mode (authentication error, not hypothesis refutation).

\textbf{External Validity}: Single language (typed Python with Pydantic), limited to HumanEval-scale programs (10--50 LOC). Generalization to Rust + Prusti or larger codebases untested.

\textbf{Construct Validity}: Speedup measured as wall-clock time (includes extraction overhead, SMT solving, dependency analysis). Extraction overhead could dominate.

\textbf{Conclusion Validity}: No statistical tests conducted (no data).
